% ============================================================
%  SOS/A FORMELSAMMLUNG — Teil 7
%  Die harmonische Balance
% ============================================================

\bereich[cKap7]{Teil 7: Harmonische Balance — Grenzzyklen}

\begin{formelreihe}
  \parbox[t]{\linewidth}{%
    {\scriptsize\scshape\color{gray!65!black} Def.: Grenzzyklus}\\[0.03em]%
    \scriptsize
    Period.\,Lösung der Zustandsgl.\\
    (geschl.\,Trajektorie); in Nähe\\
    keine weitere period.\,Lösung.\\[0.06em]
    {\scriptsize\scshape\color{gray!65!black} Arten}\\[0.03em]%
    \textbf{Stabil:} Traj.\,konvergieren\\
    von innen \& außen auf GZ\\[0.03em]
    \textbf{Instabil:} Traj.\,divergieren\\
    vom GZ weg} &
  \parbox[t]{\linewidth}{%
    {\scriptsize\scshape\color{gray!65!black} Instabile Ruhelage $\Rightarrow$}\\[0.03em]%
    \scriptsize
    – Zustand $\to$ andere Ruhelage\\[0.02em]
    – Zustand $\to\infty$\\[0.02em]
    – Zustand $\to$ \textbf{Grenzzyklus}\\
    \quad(nichtlin.\,Dauerschwingung)\\[0.02em]
    – Chaotisches Verhalten\\[0.06em]
    {\scriptsize\scshape\color{gray!65!black} NL-Standardregelkreis}\\[0.03em]%
    \scriptsize
    – NL: durch Nullpunkt, symmetr.\\
    – $G(s)$: stabiles LZI-Teilsystem\\
    – Betrachte Grundschwingung} &
  \beispielblock{Einheitskreis-GZ}{%
    \dot{x}_1\!=\!-x_2\!+\!(1\!-\!x_1^2\!-\!x_2^2)x_1\\[0.04em]
    \dot{x}_2\!=\!\phantom{-}x_1\!+\!(1\!-\!x_1^2\!-\!x_2^2)x_2\\[0.07em]
    \Rightarrow\text{GZ: Einheitskreis}\\
    \text{(stabil, ampl.-unabhäng.)}} \\[0.3em]
\end{formelreihe}
\bereichende

\bereich[cKap7]{Grafik: NL-Standardregelkreis \& Grenzzyklus-Typen}

\noindent
\begin{minipage}[c]{0.52\linewidth}\centering
\footnotesize NL-Standardregelkreis (autonom, $w{=}0$):\\[0.3em]
\begin{tikzpicture}[>=Stealth,font=\scriptsize,
    blk/.style={draw,minimum height=6.5mm,minimum width=15mm,inner sep=1pt}]
  \node (sum) [draw,circle,minimum size=4.5mm,inner sep=0pt] {};
  \draw (sum.north)--(sum.south) (sum.west)--(sum.east);
  \node[left=5.5mm of sum] (w) {$w{=}0$};
  \node[blk,right=7mm of sum,align=center] (nl) {Kennlinie\\[-0.2em]\footnotesize $K_N(v_0)$};
  \node[blk,right=8mm of nl,align=center,fill=cKap7!12] (g) {LZI\\[-0.2em]\footnotesize $G(j\omega)$};
  \coordinate[right=8mm of g] (out);
  \draw[->] (w)--(sum);
  \draw[->] (sum)--(nl) node[pos=0.42,above]{$v$};
  \draw[->] (nl)--(g) node[midway,above]{$a$}
     node[midway,below=0.3mm,cKap7!80!black,font=\tiny]{Eing.};
  \draw[->] (g)--(out) node[midway,above]{$x$}
     node[pos=0.28,below=0.3mm,cKap7!80!black,font=\tiny]{Ausg.};
  \node[above left=-0.5mm and 0.2mm of sum] {\tiny$+$};
  \node[below left=-0.5mm and 0.2mm of sum] {\tiny$-$};
  \fill (g.east)++(4mm,0) coordinate (fb) circle (0.5pt);
  \draw[->] (fb) -- ++(0,-8mm) -| (sum.south);
\end{tikzpicture}\\[0.2em]
{\scriptsize $v$ = NL-Eing.\ (Arg.\ von $K_N$);\; \textbf{LZI $G$: Eingang $a$} ($=$ NL-Ausg.), \textbf{Ausgang $x$} ($\to$ Rückf.);\; alle stat.\ Verstärker $\in G$.}
\end{minipage}\hfill
\begin{minipage}[c]{0.44\linewidth}\centering
\footnotesize Grenzzyklus-Arten (Phasenebene):\\[0.3em]
\begin{tikzpicture}[>=Stealth,font=\scriptsize,scale=0.92]
  % --- stabil: Traj. konvergieren von innen & außen auf GZ ---
  \begin{scope}
    \draw[gray,ortspfeil=0.9] plot[domain=0:250,samples=50,variable=\t]
       ({(1.28-0.55*\t/250)*cos(\t)},{(1.28-0.55*\t/250)*sin(\t)});
    \draw[gray,ortspfeil=0.9] plot[domain=0:250,samples=50,variable=\t]
       ({(0.26+0.44*\t/250)*cos(\t)},{(0.26+0.44*\t/250)*sin(\t)});
    \draw[very thick,cKap7!80!black] (0,0) circle (0.72);
    \node[cKap7!80!black] at (0,0) {\tiny GZ};
    \node at (0,-1.15) {\scriptsize\textbf{stabil} $\bigcirc$};
  \end{scope}
  % --- instabil: Traj. divergieren vom GZ weg ---
  \begin{scope}[xshift=2.9cm]
    \draw[gray,ortspfeil=0.9] plot[domain=0:250,samples=50,variable=\t]
       ({(0.75+0.53*\t/250)*cos(\t)},{(0.75+0.53*\t/250)*sin(\t)});
    \draw[gray,ortspfeil=0.9] plot[domain=0:250,samples=50,variable=\t]
       ({(0.68-0.52*\t/250)*cos(\t)},{(0.68-0.52*\t/250)*sin(\t)});
    \draw[very thick,cKap7!80!black] (0,0) circle (0.72);
    \node[cKap7!80!black] at (0,0) {\tiny GZ};
    \node at (0,-1.15) {\scriptsize\textbf{instabil} $\triangle$};
  \end{scope}
\end{tikzpicture}\\[0.2em]
{\scriptsize Traj.\ von innen \& außen: konvergent $\Rightarrow$ stabil, divergent $\Rightarrow$ instabil.}
\end{minipage}
\bereichende

\bereich[cKap7]{Beschreibungsfunktion \& Schwingungsbedingung}

\begin{formelreihe}
  \parbox[t]{\linewidth}{%
    {\scriptsize\scshape\color{gray!65!black} Grundschwingungsnäherung}\\[0.03em]%
    \scriptsize
    $v(t)\!\approx\!v_0\cos(\omega_0 t)$\\[0.03em]
    stat., symm.\,Kennlinie $\Rightarrow$\\
    Grundschw.\,$a(t)$ phasengleich\\
    mit $v(t)$\\[0.06em]
    {\scriptsize\scshape\color{gray!65!black} Beschreibungsfunktion $K_N(v_0)$}\\[0.03em]%
    $\displaystyle K_N(v_0)=\frac{a_1}{v_0}$\\[0.04em]
    {\scriptsize Fourier-Grundkoeff.\,von $a(t)$:}\\[0.02em]
    $\displaystyle a_1=\tfrac{2}{T_P}\!\int_0^{T_P}\!a(t)\cos\tfrac{2\pi t}{T_P}\,dt$\\[0.03em]
    \scriptsize NL $\to$ amplitudenabh.\,Verstärk.} &
  \parbox[t]{\linewidth}{%
    {\scriptsize\scshape\color{cKap7!80!black} Schwingungsbedingung}\\[0.03em]%
    $\displaystyle G(j\omega_0)=-\frac{1}{K_N(v_0)}$\\[0.06em]
    {\scriptsize\scshape\color{gray!65!black} Zwei-Ortskurven-Methode}\\[0.03em]%
    \scriptsize
    1.\;$G(j\omega)$-Ortskurve in kompl.\,Ebene\\[0.02em]
    2.\;Neg.\,inv.\,BF $-\tfrac{1}{K_N(v_0)}$ eintragen\\[0.02em]
    3.\;Schnittpunkt $=$ GZ-Kandidat\\[0.02em]
    \quad $v_0$: Beziff.\,der inv.-BF-Kurve\\[0.01em]
    \quad $\omega_0$: Beziff.\,der $G(j\omega)$-Kurve\\[0.01em]
    \quad\;\;bzw.\ Phasengang bei $-180^\circ$\\[0.04em]
    Periodendauer:\;$T_0=\tfrac{2\pi}{\omega_0}$\\[0.04em]
    Ampl.\,$v_0$ am Ausgang von $G(s)$} &
  \beispielblock{Schalter, $G\!=\!\tfrac{1}{(1+s)^3}$}{%
    K_N=\tfrac{4\cdot 1}{\pi v_0}\\[0.04em]
    -\tfrac{1}{K_N}=-\tfrac{\pi v_0}{4}\\[0.04em]
    \text{(neg.\,reelle Achse)}\\[0.05em]
    \Rightarrow\omega_0\!\approx\!1{,}7\\[0.02em]
    \Rightarrow T_0\!\approx\!3{,}7\\[0.02em]
    \Rightarrow v_0\!\approx\!0{,}16} \\[0.3em]
\end{formelreihe}
\bereichende

\bereich[cKap7]{Grafik: Zwei-Ortskurven-Methode (grafische GZ-Bestimmung)}

\noindent
\begin{minipage}[c]{0.53\linewidth}\centering
\begin{tikzpicture}[>=Stealth,font=\scriptsize,scale=0.92]
  % --- Achsen ---
  \draw[->] (-3.7,0)--(1.0,0) node[below]{$\operatorname{Re}$};
  \draw[->] (0,-3.2)--(0,1.0) node[left]{$\operatorname{Im}$};
  % --- Ortskurve G(jw) ---
  \draw[cKap7!80!black,thick,ortspfeil=0.5]
     plot[smooth,tension=0.7] coordinates{
       (-0.15,-2.95)(-1.5,-2.55)(-2.6,-1.6)(-2.95,-0.7)
       (-2.75,0.0)(-2.15,0.5)(-1.05,0.72)(-0.25,0.32)(0.02,0.0)};
  \node[cKap7!80!black] at (-0.75,-2.95) {$G(j\omega)$};
  \draw[cKap7!80!black,fill=white] (-1.5,-2.55) circle (1.3pt);
  \node[cKap7!80!black] at (-1.15,-2.2) {$\omega{\uparrow}$};
  % --- Neg. inv. BF (Schalter: neg. reelle Achse) ---
  \draw[red!75!black,very thick,->] (0,0)--(-3.55,0);
  \node[red!75!black,align=center] at (-1.55,0.32)
     {$-\tfrac{1}{K_N},\;v_0{\uparrow}$};
  % --- Schnittpunkt = GZ ---
  \fill (-2.75,0) circle (1.7pt);
  \node[fill=white,inner sep=1pt,align=center] at (-3.0,0.55)
     {\scriptsize\textbf{GZ}\\[-0.15em]\scriptsize$\omega_0,v_0$};
\end{tikzpicture}\\[0.1em]
{\scriptsize\textbf{Achsen:} $x{=}\operatorname{Re}\{G(j\omega)\}$,\; $y{=}\operatorname{Im}\{G(j\omega)\}$.\;
\textbf{Schnittpunkt} $G(j\omega_0){=}-\tfrac1{K_N(v_0)}$:\;
$v_0$ an $-1/K_N$-Kurve;\; $\omega_0$ an $G$-Kurve \textbf{oder} bequemer aus dem \textbf{Phasengang bei $\arg G{=}-180^\circ$};\; $T_0{=}\tfrac{2\pi}{\omega_0}$.}
\end{minipage}\hfill
\begin{minipage}[c]{0.44\linewidth}\centering
\footnotesize Stabilitätsregel am Schnittpunkt:\\[0.3em]
\begin{tikzpicture}[>=Stealth,font=\scriptsize,scale=1.0]
  % G-Kurve mit omega-Pfeil (aufwärts)
  \draw[cKap7!80!black,thick,ortspfeil=0.62]
     (0.75,-1.15) .. controls (0.18,-0.35) and (0.72,0.35) .. (0.42,1.15);
  \node[cKap7!80!black] at (1.05,1.05) {$G,\,\omega{\uparrow}$};
  % -1/K_N mit v0-Pfeil (nach links)
  \draw[red!75!black,very thick,->] (1.45,-0.1)--(-1.25,-0.1);
  \node[red!75!black] at (-0.9,0.18){$v_0{\uparrow}$};
  % Schnittpunkt
  \fill (0.42,-0.1) circle (1.5pt);
  % eps-Punkt links (stabil)
  \draw[fill=white] (-0.35,-0.1) circle (1.7pt);
  \node[align=center] at (-0.55,0.55) {\scriptsize \textbf{links}\\[-0.2em]\scriptsize$\Rightarrow$ stabil $\bigcirc$};
  \node[align=center] at (1.15,-0.62) {\scriptsize \textbf{rechts}\\[-0.2em]\scriptsize$\Rightarrow$ instab.\ $\triangle$};
\end{tikzpicture}\\[0.1em]
{\scriptsize $v_0$ um $\varepsilon$ erhöhen: liegt $-\tfrac1{K_N}(v_0{+}\varepsilon)$
\textbf{links} der $G$-Kurve (in $\omega$-Richtung) $\Rightarrow$ \textbf{stabiler} GZ.}
\end{minipage}
\bereichende

\bereich[cKap7]{Rezept: Dauerschwingung (GZ) bestimmen \& Amplituden}

\begin{rezept}
\rs{\textbf{Standardregelkreis bilden:} alle LZI-Blöcke zu
  $G(s)=V_1 V_2\,F(s)$ zusammenfassen (statische Verstärker $\cdot$ dynam.\ Teil), NL als Block $K_N(v_0)$ separat.}
\rs{$G(j\omega)$-Ortskurve (beziffert mit $\omega$) in kompl.\ Ebene eintragen.}
\rs{Neg.\ inv.\ BF $-\tfrac{1}{K_N(v_0)}$ eintragen (Lage s.\ BF-Tabelle: reelle Achse / Hysterese-Parallele).}
\rs{\textbf{Schwingungsbed.}\ $G(j\omega_0)=-\tfrac{1}{K_N(v_0)}$ = \textbf{Schnittpunkt} beider Kurven:\;
  $v_0$ an der $-1/K_N$-Kurve ablesen;\;
  $\omega_0$ an der $G$-Kurve \textbf{oder} — bei Schnitt auf der neg.\ reellen Achse ($\arg G{=}-180^\circ$) — bequemer im \textbf{Phasengang beim Durchgang durch $-180^\circ$} ablesen.}
\rs{Periodendauer $T_0=\tfrac{2\pi}{\omega_0}$.}
\rs{\textbf{Stabilität:} $v_0$ leicht erhöhen ($+\varepsilon$). Liegt $-\tfrac{1}{K_N(v_0+\varepsilon)}$ \textbf{links} der $G(j\omega)$-Kurve (Durchlaufsinn wachsender $\omega$)?\; Ja $\Rightarrow$ stabiler GZ, Nein $\Rightarrow$ instabil.}
\end{rezept}

\miniexample{Zweipunkt-Schalter (Höhe $h$), $V_1\!=\!V_2\!=\!1$}{%
$F(j\omega)$ schneide die neg.\ reelle Achse bei $-1$; dort sei $\omega_0=0{,}65$ abgelesen.
Schalter: $K_N=\tfrac{4h}{\pi v_0}$, also $-\tfrac{1}{K_N}=-\tfrac{\pi v_0}{4h}$ (reell).
Schnittpunkt: $-\tfrac{\pi v_0}{4h}=-1\;\Rightarrow\;v_0=\tfrac{4h}{\pi}$.\quad
$T_0=\tfrac{2\pi}{0{,}65}\approx 9{,}7$. Schalter $\Rightarrow$ stets stabiler GZ.}

\bereichende

\bereich[cKap7]{Rezept: Signal-Amplituden im Regelkreis (Propagation)}

\noindent\footnotesize $v_0$ (HB) $=$ Amplitude am \textbf{NL-Eingang} (Arg.\ von $K_N$).\;
\textbf{NL-Ausgang} (Grundschw.): $a_1=K_N(v_0)\,v_0$ — \textbf{Schalter:} $a_1=\tfrac{4h}{\pi}=$ \emph{konst.}, unabh.\ von $v_0$ (\& von $\alpha$)!\;
Übrige Amplituden per \textbf{Zurück-/Vorwärtsrechnen} durch die Blöcke (je Grundschw., Frequenz $\omega_0$):\par\vspace{0.15em}
\begin{rezept}
\rs{Statischer Verstärker $V$:\; Ausgangsamplitude $=V\cdot$ Eingangsamplitude.}
\rs{Dynam.\ Block $F$:\; Ausgangsamplitude $=|F(j\omega_0)|\cdot$ Eingangsamplitude.\;
  $|F(j\omega_0)|$ aus Ortskurve (Betrag des Zeigers) bzw.\ $=\tfrac{|G(j\omega_0)|}{V_1V_2}$.}
\rs{Kette $b\!\to\!V_1\!\to\!u\!\to\!F\!\to\!x\!\to\!V_2\!\to\!v$:\;
  $x_0=\tfrac{v_0}{V_2}$,\;\; $u_0=\tfrac{x_0}{|F(j\omega_0)|}$,\;\; $b_0=\tfrac{u_0}{V_1}$.}
\end{rezept}

\miniexample{$V_1\!=\!V_2\!=\!2$, abgelesen $v_0=5$, $|F(j\omega_0)|=1$}{%
$x_0=\tfrac{v_0}{V_2}=2{,}5$;\quad $u_0=\tfrac{x_0}{|F(j\omega_0)|}=2{,}5$;\quad $b_0=\tfrac{u_0}{V_1}=1{,}25$.}

\bereichende

\bereich[cKap7]{Rezept: Entwurf — Verstärkung / Schalthöhe für Soll-Amplitude}

\begin{rezept}
\rs{Soll-Amplitude (z.\,B.\ $x_0$) über die Blockkette in $v_0$ am NL-Eingang umrechnen (s.\ Propagation).}
\rs{$K_N(v_0)$ aus BF bestimmen (Graph ablesen oder Formel einsetzen).}
\rs{Schwingungsbed.\ am Schnittpunkt erzwingen: freie Verstärkung so wählen, dass $G(j\omega_0)=-\tfrac{1}{K_N(v_0)}$.
  Schneidet $F$ die neg.\ reelle Achse beim Wert $r$, gilt $V_1V_2\,r=-\tfrac{1}{K_N(v_0)}$ $\Rightarrow$ freie Verstärkung lösen.}
\end{rezept}

\miniexample{Dreipunkt-Schalter ($h\!=\!d\!=\!1$), $V_2\!=\!1$: welches $V_1$ für $x_0=4$?}{%
$V_2=1\Rightarrow v_0=x_0=4$. BF-Graph: $K_N(4)=0{,}3$ $\Rightarrow -\tfrac{1}{K_N}=-3{,}33$.
$F$ schneidet neg.\ reelle Achse bei $-1$ $\Rightarrow V_1\cdot(-1)=-3{,}33$ $\Rightarrow \boxed{V_1=3{,}33}$.}

\fallstrick{%
\textbullet\ \textbf{Drei Amplituden unterscheiden:} $v_0$ = NL-\emph{Eingang}; $a_1{=}K_N v_0$ = NL-\emph{Ausgang} (Schalter $\tfrac{4h}{\pi}$, konst.); dahinter $\cdot|G(j\omega_0)|$ = $G$-Ausgang. \emph{Notation:} FS-$a$ = NL-Ausgang ($\hat{=}$ FS-$x$)!\quad
\textbullet\ \textbf{Zweipunkt-Schalter:} neg.-inv.\ BF deckt die \emph{ganze} neg.\ reelle Achse ab $\Rightarrow$ $\omega_0$ (und $T_0$) liegt fest durch $G$; $h$ ändert nur $v_0$, \textbf{nicht} $T_0$.\quad
\textbullet\ \textbf{Schalter $\Rightarrow$ Vorfaktor $\alpha$ wirkungslos auf $a_0$:} Schalter-Ausgang $a_1{=}\tfrac{4h}{\pi}$ hängt \emph{nicht} von der Eingangsamplitude ab, und $\alpha$ (reell) verschiebt den $-180^\circ$-Durchgang nicht $\Rightarrow$ $a_0{=}|G(j\omega_0)|\tfrac{4h}{\pi}$ liegt \textbf{fest}; $\alpha$ skaliert nur $v_0{=}\alpha a_0$. Soll-$a_0$ \textbf{nicht} über $\alpha$ einstellbar! (\textbf{Begrenzer:} $K_N$ hängt von $v_0$ ab $\Rightarrow$ $\alpha$ \emph{wirkt}, GZ ein-/ausschaltbar.)\quad
\textbullet\ \textbf{Dreipunkt/Begrenzer/Totzone:} $K_N=0$ für $v_0<d$; neg.-inv.\ BF beginnt erst bei endlichem Wert $\Rightarrow$ evtl.\ \textbf{kein Schnittpunkt} $\Rightarrow$ HB sagt keine Dauerschwingung.\quad
\textbullet\ \textbf{Kein Schnittpunkt $\not\Rightarrow$ kein GZ!} HB ist Näherung, exakt nur bei ausgeprägter TP-Wirkung von $G$ (Oberwellen gedämpft).\quad
\textbullet\ \textbf{Hysterese:} $K_N\in\mathbb{C}$ $\Rightarrow$ neg.-inv.\ BF \emph{nicht} auf der reellen Achse ($\operatorname{Im}=-\tfrac{d\pi}{4h}=$ const).}

\bereichende

\needspace{17\baselineskip}
\bereich[cKap7]{Beschreibungsfunktionen — Übersicht}

\noindent
\begin{minipage}[c]{0.45\linewidth}\centering
\begin{tikzpicture}[>=Stealth,font=\scriptsize,scale=1.0]
  \draw[->] (-4.1,0)--(0.8,0) node[below]{$\operatorname{Re}$};
  \draw[->] (0,-1.25)--(0,0.7) node[left]{$\operatorname{Im}$};
  % Schalter: ganze neg. reelle Achse
  \draw[red!75!black,very thick,->] (0,0)--(-3.95,0);
  % Dreipunkt/Begrenzer/Totzone: ab endlichem Wert (leicht versetzt)
  \draw[orange!85!black,very thick,->] (-1.35,0.16)--(-3.95,0.16);
  \fill[orange!85!black] (-1.35,0.16) circle (1.2pt);
  % Schalter mit Hysterese: Parallele zur reellen Achse
  \draw[blue!70!black,very thick,->] (0,-0.85)--(-3.95,-0.85);
  \node[blue!70!black,right] at (0.05,-0.85) {$-\tfrac{d\pi}{4h}$};
\end{tikzpicture}
\end{minipage}\hfill
\begin{minipage}[c]{0.53\linewidth}\footnotesize
{\color{red!75!black}\rule{9pt}{2.4pt}}\; \textbf{Schalter:} ganze neg.\ reelle Achse ($0\!\to\!-\infty$).\\[0.25em]
{\color{orange!85!black}\rule{9pt}{2.4pt}}\; \textbf{Dreipkt./Begr./Totzone:} reelle Achse \emph{ab endl.\ Wert} ($K_N{=}0$ für $v_0{<}d$).\\[0.25em]
{\color{blue!70!black}\rule{9pt}{2.4pt}}\; \textbf{Schalter m.\ Hyst.:} Parallele zur reellen Achse, $\operatorname{Im}{=}-\tfrac{d\pi}{4h}{=}$const.\\[0.25em]
\textit{ohne Gedächtnis $\Rightarrow K_N\in\mathbb{R}$ (reelle Achse); mit Gedächtnis $\Rightarrow K_N\in\mathbb{C}$.}\\[0.25em]
\textbf{Achsen} der Ortskurve: $x{=}\operatorname{Re}\{G(j\omega)\}$, $y{=}\operatorname{Im}\{G(j\omega)\}$ Bei gedächtnisloser NL liegt $-1/K_N$ auf der $x$-Achse (neg.\ reell) $\Rightarrow$ Schnittpunkt bei $\operatorname{Im}\{G\}{=}0$, dort $\arg G{=}-180^\circ$.
\end{minipage}\par\vspace{0.2em}
\noindent\footnotesize\setlength{\tabcolsep}{2pt}%
{\renewcommand{\arraystretch}{1.15}%
% --- Mini-Kennlinien (einheitlicher Stil) ---
\newcommand{\knlaxes}{\draw[->,gray!55](-1.5,0)--(1.5,0);\draw[->,gray!55](0,-1.25)--(0,1.5);}
\begin{tabular}{@{}>{\centering\arraybackslash}m{0.19\linewidth}@{\hspace{2pt}}m{0.48\linewidth}@{\hspace{3pt}}m{0.29\linewidth}@{}}
  \scriptsize\scshape\color{gray!65!black}Kennlinie \& NL &
  \scriptsize\scshape\color{gray!65!black}Beschreibungsfunktion $K_N(v_0)$ &
  \scriptsize\scshape\color{gray!65!black}Neg.\,inv.\,BF\;$-1/K_N$ \\[0.1em]
  % --- Schalter (Zweipunkt) ---
  \begin{tikzpicture}[scale=0.34,>=Stealth,line width=0.6pt]
    \knlaxes
    \draw[very thick](-1.3,-0.9)--(0,-0.9)--(0,0.9)--(1.3,0.9);
    \node[font=\tiny] at (-0.32,1.18){$h$};
    \node[font=\footnotesize,anchor=north] at (0,-1.55){Schalter};
    \node[font=\footnotesize,anchor=south,opacity=0] at (0,1.55){Schalter};
  \end{tikzpicture} &
  $\tfrac{4h}{\pi v_0}$ &
  $-\tfrac{\pi v_0}{4h}$ (reell)\\[0.2em]
  % --- Schalter mit Hysterese ---
  \begin{tikzpicture}[scale=0.34,>=Stealth,line width=0.6pt]
    \knlaxes
    \draw[very thick](-1.3,-0.9)--(0.55,-0.9)--(0.55,0.9)--(1.3,0.9);
    \draw[very thick](1.3,0.9)--(-0.55,0.9)--(-0.55,-0.9)--(-1.3,-0.9);
    \draw[very thick,->](-0.45,-0.9)--(-0.15,-0.9);
    \draw[very thick,->](0.45,0.9)--(0.15,0.9);
    \node[font=\tiny] at (-0.32,1.18){$h$};
    \node[font=\tiny] at (0.78,-1.12){$d$};
    \node[font=\footnotesize,anchor=north] at (0,-1.55){Schalter m.\,Hyst.};
    \node[font=\footnotesize,anchor=south,opacity=0] at (0,1.55){Schalter m.\,Hyst.};
  \end{tikzpicture} &
  \begin{minipage}[c]{\linewidth}
    \phantom{\scriptsize$v_0$}\\[0.12em]
    $\tfrac{4h}{\pi v_0}\sqrt{1\!-\!\tfrac{d^2}{v_0^2}}-j\tfrac{4dh}{\pi v_0^2}$\\[0.12em]
    {\scriptsize\itshape $v_0\!\geq\!d$ (sonst kommt $v_0\!<\!d$ nicht vor)}
  \end{minipage} &
  \begin{minipage}[c]{\linewidth}
    \phantom{\scriptsize$v_0$}\\[0.12em]
    $-\tfrac{\pi v_0}{4h}\sqrt{1\!-\!\tfrac{d^2}{v_0^2}}-j\tfrac{d\pi}{4h}$\\[0.12em]
    {\scriptsize\itshape $\operatorname{Im}=-\tfrac{d\pi}{4h}=\mathrm{konst.}$}
  \end{minipage}\\[0.2em]
  % --- Dreipunktschalter ---
  \begin{tikzpicture}[scale=0.34,>=Stealth,line width=0.6pt]
    \knlaxes
    \draw[very thick](-1.3,-0.9)--(-0.45,-0.9)--(-0.45,0)--(0.45,0)--(0.45,0.9)--(1.3,0.9);
    \node[font=\tiny] at (-0.32,1.18){$h$};
    \node[font=\tiny] at (0.62,-0.26){$d$};
    \node[font=\footnotesize,anchor=north] at (0,-1.55){Dreipunktschalter};
    \node[font=\footnotesize,anchor=south,opacity=0] at (0,1.55){Dreipunktschalter};
  \end{tikzpicture} &
  $\begin{cases}0 & v_0\!<\!d\\\tfrac{4h\sqrt{v_0^2-d^2}}{\pi v_0^2} & v_0\!\geq\!d\end{cases}$ &
  reell ($v_0\!\geq\!d$)\\[0.5em]
  % --- Begrenzer (Sättigung) ---
  \begin{tikzpicture}[scale=0.34,>=Stealth,line width=0.6pt]
    \knlaxes
    \draw[very thick](-1.3,-0.9)--(-0.55,-0.9)--(0.55,0.9)--(1.3,0.9);
    \node[font=\tiny] at (-0.32,1.18){$h$};
    \node[font=\tiny] at (0.72,0.22){$d$};
    \node[font=\footnotesize,anchor=north] at (0,-1.55){Begrenzer};
    \node[font=\footnotesize,anchor=south,opacity=0] at (0,1.55){Begrenzer};
  \end{tikzpicture} &
  $\begin{cases}\tfrac{h}{d} & v_0\!<\!d\\[0.1em]\tfrac{2h}{d\pi}\!\left(\arcsin\tfrac{d}{v_0}\!+\!\tfrac{d}{v_0}\sqrt{1\!-\!\tfrac{d^2}{v_0^2}}\right) & v_0\!\geq\!d\end{cases}$ &
  reell\\[0.65em]
  % --- Totzone ---
  \begin{tikzpicture}[scale=0.34,>=Stealth,line width=0.6pt]
    \knlaxes
    \draw[very thick](-1.3,-0.85)--(-0.45,0)--(0.45,0)--(1.3,0.85);
    \node[font=\tiny] at (1.12,1.15){$m$};
    \node[font=\tiny] at (0.6,-0.26){$d$};
    \node[font=\footnotesize,anchor=north] at (0,-1.55){Totzone};
    \node[font=\footnotesize,anchor=south,opacity=0] at (0,1.55){Totzone};
  \end{tikzpicture} &
  $\begin{cases}0 & v_0\!<\!d\\[0.1em]m\!-\!\tfrac{2m}{\pi}\!\left(\arcsin\tfrac{d}{v_0}\!+\!\tfrac{d}{v_0}\sqrt{1\!-\!\tfrac{d^2}{v_0^2}}\right) & v_0\!\geq\!d\end{cases}$ &
  reell\\
\end{tabular}}\par
\vspace{0.15em}
\bereichende

\bereich[cKap7]{Stabilität von Grenzzyklen}

\begin{formelreihe}
  \parbox[t]{\linewidth}{%
    {\scriptsize\scshape\color{gray!65!black} Stabilitätsprüfung am Schnittpunkt}\\[0.03em]%
    \scriptsize
    Am Schnittpunkt gilt $G(j\omega_0)=-\tfrac{1}{K_N(v_0)}$\\[0.04em]
    Erhöhe $v_0$ leicht ($v_0\!+\!\varepsilon$):\\[0.03em]
    Liegt $-\tfrac{1}{K_N(v_0\!+\!\varepsilon)}$\\
    \quad\textbf{links} der $G(j\omega)$-Ortskurve?\\[0.04em]
    $\Rightarrow$ \textbf{Ja}: stabiler GZ\;$\bigcirc$\\[0.03em]
    $\Rightarrow$ \textbf{Nein}: instabiler GZ\;$\triangle$\\[0.05em]
    \textit{\color{red!70!black} Stabil: System zeigt nach dem Einschalten eine nichtlin.\ Dauerschwingung.}} &
  \parbox[t]{\linewidth}{%
    {\scriptsize\scshape\color{gray!65!black} Gültigkeit der Methode}\\[0.03em]%
    \scriptsize
    exakt nur bei ausgepr.\,TP-Wirkung\\
    von $G(s)$ (Oberwellen gedämpft)\\[0.04em]
    \textbf{Kein Schnittpunkt} $\not\Rightarrow$ kein GZ!\\[0.03em]
    Bsp.: PT1\,+\,Schalter\,m.\,Hyst.\\
    $\Rightarrow$ GZ trotzdem\\[0.04em]
    ohne Schnittpunkt: keine\\
    Aussage über Schwingung.} &
  \beispielblock{Vorgehen kompakt}{%
    \text{1. }G(j\omega)\text{-Ortskurve}\\[0.03em]
    \text{2. }-\tfrac{1}{K_N(v_0)}\text{-Kurve}\\[0.03em]
    \text{3. Schnittpkt.\!}\to\!v_0,\omega_0\\[0.03em]
    \text{4. }T_0=\tfrac{2\pi}{\omega_0}\\[0.04em]
    \text{5. }v_0{\uparrow}\!:\text{ links?}\Rightarrow\!\text{stab.}} \\[0.3em]
\end{formelreihe}
\bereichende

\bereich[cKap7]{Sonderfall: mehrere Schnittpunkte (schneller \& langsamer GZ)}

\noindent
\begin{minipage}[c]{0.52\linewidth}\centering
\begin{tikzpicture}[>=Stealth,font=\scriptsize,scale=0.92]
  \draw[->] (-3.1,0)--(0.7,0) node[below]{$\operatorname{Re}$};
  \draw[->] (0,-2.45)--(0,1.15) node[left]{$\operatorname{Im}$};
  % -1/K_N = ganze neg. reelle Achse (Zweipunktschalter, Hyst -> 0)
  \draw[red!75!black,very thick,->] (0,0)--(-3.0,0);
  % Ortskurve G(jw): schneidet neg. reelle Achse zweimal
  \draw[cKap7!80!black,thick,ortspfeil=0.45]
     plot[smooth,tension=0.75] coordinates{
       (-2.0,0.85)(-2.5,0.0)(-2.78,-1.15)(-1.8,-2.05)
       (-0.7,-1.55)(-0.2,-0.55)(-0.13,0.0)(0.0,0.03)};
  \node[cKap7!80!black] at (-1.5,-2.15) {$G(j\omega),\,\omega{\uparrow}$};
  % Schnittpunkt B (weit außen, langsam, instabil)
  \fill (-2.5,0) circle (1.7pt);
  \node[align=center,fill=white,inner sep=1pt] at (-2.5,0.5)
     {\scriptsize$\triangle$ instabil\\[-0.15em]\scriptsize langsam ($\omega{\downarrow}$)};
  % Detail-Box um Null
  \draw[gray,dashed] (-0.45,-0.4) rectangle (0.2,0.4);
  \node[gray] at (-0.15,-0.62) {\scriptsize Detail};
\end{tikzpicture}
\end{minipage}\hfill
\begin{minipage}[c]{0.44\linewidth}\centering
\footnotesize Detail nahe $0$ ($\times 10^{-3}$):\\[0.2em]
\begin{tikzpicture}[>=Stealth,font=\scriptsize,scale=1.0]
  \draw[->] (-2.2,0)--(0.6,0) node[below]{$\operatorname{Re}$};
  \draw[->] (0,-0.95)--(0,0.85) node[left]{$\operatorname{Im}$};
  \draw[red!75!black,very thick,->] (0,0)--(-2.1,0);
  \draw[cKap7!80!black,thick,ortspfeil=0.55]
     plot[smooth,tension=0.7] coordinates{
       (-1.95,0.55)(-1.4,0.0)(-0.8,-0.42)(-0.2,-0.2)(0.0,0.0)};
  \node[cKap7!80!black] at (-1.75,0.6) {$G,\,\omega{\uparrow}$};
  % Schnittpunkt A (nahe 0, schnell, stabil)
  \fill (-1.4,0) circle (1.7pt);
  \node[align=center,fill=white,inner sep=1pt] at (-0.75,0.5)
     {\scriptsize$\bigcirc$ stabil\\[-0.15em]\scriptsize schnell ($\omega{\uparrow}$)};
\end{tikzpicture}
\end{minipage}\par\vspace{0.25em}
\noindent\footnotesize
\textbf{Zweipunktschalter, Hyst.\ $\to 0$} $\Rightarrow -\tfrac1{K_N}$ = \emph{ganze} neg.\ reelle Achse.
Ortskurve schneidet sie \textbf{zweimal} $\Rightarrow$ \textbf{zwei GZ}. Links/Rechts-Regel je Schnittpunkt:
\textbullet\ \textbf{schneller} GZ nahe $0$ (großes $\omega$, $T_0{\to}0$): \textbf{stabil} $\bigcirc$ — stellt sich \emph{ohne} Anregung ein.\quad
\textbullet\ \textbf{langsamer} GZ weit außen (kleines $\omega$, großes $T_0$): \textbf{instabil} $\triangle$ — nur mit Anfangs-Anregung erreichbar (Schwelle zwischen den Einzugsgebieten).
\bereichende
